\section{Conclusion}
\label{sec:conclusion}

We formulated complex multi-hop retrieval as \emph{declarative evidence
execution}: a question is compiled into typed, dependent evidence requirements;
a deterministic importance-weighted allocator assigns legal orders, physical
implementations, and budget allocations to those requirements under a matched
budget; and a deterministic executor acquires evidence with hard budget
enforcement and full provenance. We introduced a typed evidence algebra whose
state-transition semantics (\code{apply\_observation}) make requirement
progress first-class and auditable, and we derived a chain law pinning
requirement importance to dependency depth ($\tau = 2d-1$) on well-defined
chains, so the estimator is construction-derived rather than learned.

The evaluation does \emph{not} establish a universal accuracy or cost win, and
we are deliberate about that boundary. On the frozen SEALED\_TEST set (8,661
questions; model \code{qwen3.5-9b}; matched retrieval budget $=8$) the
chain-rule allocator is significantly \emph{worse} than static on
2WikiMultiHopQA (EM $47.8\%$ vs.\ $51.7\%$, paired bootstrap $\Delta$EM
$-3.9$pt CI$[-5.0,-2.9]$, McNemar $p<0.0001$), statistically tied on HotpotQA
(flat $55.4\%$ / static $54.1\%$ / chain $54.2\%$; chain-vs-static McNemar
$p{=}0.80$), and only marginally best on MuSiQue (chain $27.5\%$, McNemar
$p{=}0.064$). The flat arm---which performs no dependency-tree search---is
strongest or tied on every dataset. The matched-budget frontier is similarly
heterogeneous: mean retrieval calls are near-identical between chain and static
(HotpotQA $-0.20$, 2Wiki $+0.01$, MuSiQue $-0.09$ calls/question), while the
chain arm spends more LLM calls on HotpotQA ($+24.0$ per question,
$p<0.0001$). The honest statement is that allocator benefit is conditional on
dependency depth: it helps on the deep chains of MuSiQue and is
neutral-to-harmful on the shallow plans that dominate 2WikiMultiHopQA, where
the chain law leaves little allocation freedom. We report the absence of a
global win with the same rigor as a win would demand.
We treat runtime re-optimization over alternative physical plans as future work
under a branching execution model, because on the chain-only topology this
system runs today it is vacuous.

More broadly, this study suggests that complex retrieval can be treated as an
execution problem in which evidence goals remain stable while physical
acquisition strategies adapt to observed state and resource constraints. The
honest path forward is to make the evidence requirements and their satisfaction
state the interface between declarative programs and physical optimization, and
to evaluate every claim---including the absence of a benefit---against a
matched budget with the same rigor as the presence of one.

\paragraph{Limitations}
The matched-budget headline numbers rest on the frozen SEALED\_TEST set
($n{=}2{,}863/4{,}931/867$ HotpotQA / 2WikiMultiHop / MuSiQue, pooled
$n{=}8{,}661$) under one decoder (\code{qwen3.5-9b} at temperature $0$) through
a shared provider; generalizability to other decoders is untested. Of the
$25{,}983$ executions, $24{,}067$ completed ($92.6\%$); on the solvable subset
($25{,}201$, excluding $22$ residual infrastructure failures and $760$
deterministic plan--question mismatches the compiler correctly rejects) the
completed rate is $95.5\%$, and $1{,}134$ questions exceeded the matched
retrieval budget. The budget hits are not uniformly chain-skewed: they
concentrate on the deep HotpotQA chain arm ($323$) and on the shallow 2Wiki
static/flat arms ($265$/$246$), whereas the 2Wiki chain arm hits the ceiling
only twice ($2$).
Evidence corpora are the public HotpotQA/2Wiki/MuSiQue passage indices; the
heterogeneous-evidence result (HoVer/FEVEROUS pilot) is a small transfer
proof-of-concept, not a Wikipedia-scale retrieval benchmark. The chain law
$\tau = 2d - 1$ is exact on well-defined \emph{chains} where the compiled
requirement order respects topological dependencies; on general (branching)
graphs the enumerate-index proxy is a heuristic and is outside the claim's
scope. Cost accounting covers realized retrieval and LLM calls under the
matched budget; latency and monetary cost are reported for the optimizer
overhead but are not the primary comparison axis.
